\documentclass[12pt, b5paper]{article}
\usepackage{ctex}

\usepackage{geometry}
\geometry{b5paper,left=2cm,right=1cm,top=2.5cm,bottom=2.5cm}
\usepackage{chemformula} %化学公式宏包
\usepackage[version=4]{mhchem} %化学公式宏包
\usepackage{siunitx} %数字，单位宏包
\sisetup{
	separate-uncertainty = true,
	inter-unit-product = \ensuremath{{}\cdot{}}
} %单位间加小圆点
\usepackage{tikz} %Tikz绘图包
\usetikzlibrary{cd}
\usetikzlibrary{chains}
\usetikzlibrary{arrows}
\usetikzlibrary{arrows.meta} %画箭头
\usetikzlibrary{commutative-diagrams} %交换图
\usetikzlibrary{positioning,automata}
\usepackage[all]{xy}
\usepackage{booktabs} %三线表
\usepackage{multirow} %合并单元格
\usepackage{makecell} %表格内强制换行
\usepackage{array}
\usepackage{booktabs} %控制表格行高
\usepackage{colortbl}  %彩色表格需要加载的宏包
\usepackage{amsmath}
\usepackage{unicode-math}
\usepackage{fontspec} %字体
\newcommand*{\cred}{\textcolor{red}}

\setmainfont{Times New Roman}
\setmathfont{XITS Math}

\setlength{\parindent}{0pt} %无缩进
\usepackage{setspace}
\usepackage{fancyhdr} %设置页码
\usepackage{lastpage} %获取总页数
\pagestyle{fancy} %fancyhdr宏包新增的页面风格
\renewcommand\headrulewidth{0pt} %取消页眉中的装饰分割线
\fancyhf{}
\cfoot{\footnotesize 第 \thepage\ 页(共 \pageref{LastPage}页)}%当前页 of 总页数

\begin{document}
\begin{spacing}{1.625}

\centerline{{\huge 河北农业大学课程考试试卷}}

\centerline{\large 2023-2024学年第\underline{\ 1\ }学期}

考试科目：\underline{化学反应工程} \hfill 姓名：\underline{ \qquad \qquad} \hfill 学号：\underline{ \qquad \qquad \qquad \qquad}

学院：\underline{理工学院} \hfill 专业班级：\underline{\qquad \qquad \qquad \qquad} \hfill 卷别：\cred{B} \hfill 考核方式：\cred{开卷}

（注：考生务必将答案写在答题纸上，标清楚大小题号，写在本试卷上无效）

\bigskip

{\bfseries 一、简答题（共40分）}

1. (10分) 对于如下平行反应
\[
	\ch{A + B -> P}
\]
\[
	\ch{A + B -> Q}
\]
设P为目的产物且第一个反应的活化能小于第二个反应的活化能，应选择什么样的操作温度进行生产？

\bigskip
\bigskip

2. (10分) 从停留时间分布的角度解释活塞流反应器和全混流反应器反应效果不同的原因。

\bigskip
\bigskip

3. (10分) 气体在多孔介质中的扩散有哪些形式？如何判定？

\bigskip
\bigskip

4. (10分) 试叙述气固相催化反应中梯尔模数$\phi$和有效扩散系数$\eta$的物理意义。

\newpage
{\bfseries 二、计算题（共60分）}

1. (20分) 某液相均相反应\ch{A -> P}为一级不可逆反应，当反应温度100℃时，反应速率常数为\SI{1.5}{h^{-1}}，进料中A的初始浓度为2 mol/L。要求P的产量为270 mol/h，最终转化率为90\%。求下列几种情况下的反应体积：

（1）反应在间歇釜式反应器中进行，辅助时间为0.5 h；

（2）反应在全混流釜式反应器中进行；

（3）反应在活塞流管式反应器中进行。

\bigskip
\bigskip

2. (20分) 在充填直径为9 mm、高为7 mm的圆柱形铁铬催化剂的固定床反应器中，0.6865 MPa下进行水煤气变换反应。反应气体的平均相对分子质量为18.96，质量速度（按空床计算）为\SI{0.936}{kg/(s.m^2)}。设床层的平均温度为689 K，反应气体的黏度等于\SI{2.5e-5}{Pa.s}。已知催化剂的颗粒密度和床层的堆密度分别为\SI{2000}{kg/m^3}及\SI{1400}{kg/m^3}。试计算单位床层高度的压力降。

\bigskip
\bigskip

3. (20分) 在直径8 mm的薄片催化剂上进行丁烷脱氢反应，其反应速率方程为
\vspace{-1.5em}
\[
	r_{\mathrm{A}} =k_{\mathrm{W}} c_{\mathrm{A}} \quad \mathrm{mol} /(\mathrm{g} \cdot \mathrm{s})
	\vspace{-1.5em}
\]
0.1013 MPa及773 K时，$k_\mathrm{W} = \SI{0.92}{cm^3/(g.s)}$。若催化剂平均孔半径为\SI{4e-7}{cm}，孔容为\SI{0.35}{cm^3/g}，曲节因子等于2.5。试计算内扩散有效因子。

（注：$\tanh (x) = \frac{\mathrm{e}^x - \mathrm{e}^{-x}}{\mathrm{e}^x + \mathrm{e}^{-x}}$。当$x>3$时，$\tanh (x) = 1$）


\end{spacing}
\end{document}
